\begin{onehalfspace}

Questa sezione conserva la trattazione estesa e i protocolli dei laboratori
introduttivi sintetizzati nella Sezione~\ref{sec:repetition} e nella
Sezione~\ref{sec:steane} del Capitolo~\ref{chap:risultati}. Le curve e i
risultati decisivi restano nel corpo; qui sono raccolti gli strumenti
concettuali, le lookup table e le assunzioni necessarie a verificarli.

\subsection{Dalla ridondanza alla codifica per stabilizzatori}
\label{app:qec_richiami}

L'idea classica di trasmettere tre copie di un bit e decidere a maggioranza
non si trasferisce direttamente al caso quantistico. Il teorema di
no-cloning~\cite{wootters1982} impedisce di copiare uno stato ignoto e la
misura diretta distruggerebbe la sovrapposizione da proteggere. La
\ac{QEC} codifica invece l'informazione di un qubit logico nelle
correlazioni di più qubit fisici~\cite{shorcorrection1995,steane1996,roffe2019}.
Operatori ausiliari rivelano quale classe di errore si è verificata senza
rivelare i coefficienti dello stato logico.

Il ciclo generale comprende: (i) codifica di $\ket{\psi_L}$; (ii) azione
del rumore; (iii) estrazione della sindrome mediante ancille; (iv)
decodifica classica; (v) correzione fisica o aggiornamento del \emph{Pauli
frame}. Uno stabilizzatore $S$ lascia invariati gli stati del codice,
$S\ket{\psi_L}=\ket{\psi_L}$: un esito $-1$ segnala l'uscita dal
sottospazio valido senza distinguere gli stati logici che lo compongono.

Un codice $[[n,k,d]]$ usa $n$ qubit fisici per $k$ qubit logici e ha
distanza $d$, pari al peso minimo di un errore logico non rilevato. Può
correggere
\begin{equation}
 t=\left\lfloor\frac{d-1}{2}\right\rfloor
 \label{eq:distanza_t}
\end{equation}
errori arbitrari. Il teorema della soglia garantisce in linea di principio
un regime in cui l'aumento della distanza riduce l'errore
logico~\cite{aharonov1997fault}; la sua osservazione quantitativa è oggetto
della Sezione~\ref{sec:surface_risultati}.

\begin{table}[H]
\centering
\scriptsize
\begin{tabularx}{\textwidth}{@{}p{2.0cm}p{1.6cm}XXX@{}}
\toprule
\textbf{Codice} & \textbf{Parametri} & \textbf{Corregge} & \textbf{Pro / Contro} & \textbf{Uso nella tesi} \\
\midrule
Repetition (bit-flip) & $[[3,1,1]]$ su X & solo bit-flip & semplice / non corregge phase-flip & primo laboratorio \\
Repetition (phase-flip) & 3 qubit, base H & solo phase-flip & mostra dualità X/Z / parziale & controllo duale \\
Shor code & $[[9,1,3]]$ & errore arbitrario singolo & storico, intuitivo / più pesante & richiamo didattico~\cite{shorcorrection1995} \\
Steane code & $[[7,1,3]]$ & errore arbitrario singolo & \ac{CSS}, compatto / distanza bassa & secondo laboratorio~\cite{steane1996} \\
Five-qubit & $[[5,1,3]]$ & errore arbitrario singolo & minimo / non \ac{CSS} & confronto teorico \\
Surface code & $d^2$ dati $+O(d^2)$ ancille & errori locali su griglia 2D & scalabile / overhead elevato & campagna principale~\cite{fowler2012} \\
\bottomrule
\end{tabularx}
\caption[Confronto tra codici correttori quantistici]{Confronto esteso e
progressione metodologica: repetition per introdurre la sindrome, Steane
per un errore arbitrario, surface code per soglia e scalabilità.}
\label{tab:confronto_codici}
\end{table}

\subsection{Protocollo completo del codice a ripetizione}
\label{app:qec_repetition}

Il codice contro errori X codifica $\ket{0_L}=\ket{000}$ e
$\ket{1_L}=\ket{111}$: lo stato
$\alpha\ket{0}+\beta\ket{1}$ diventa
$\alpha\ket{000}+\beta\ket{111}$, uno stato entangled e non tre copie.
Gli stabilizzatori $Z_1Z_2$ e $Z_2Z_3$ misurano la parità tramite due
ancille. Il duale nella base di Hadamard usa $X_1X_2$ e $X_2X_3$ e
protegge dagli errori Z; la concatenazione delle due protezioni conduce
al codice di Shor a nove qubit.

Il primo controllo inietta deterministicamente un errore su ciascun qubit
dato e verifica che la sindrome identifichi la posizione senza misurare
l'informazione logica.

\begin{table}[H]
\centering
\begin{tabular}{lccc}
\toprule
\textbf{Errore iniettato} & \textbf{Sindrome $(a_0,a_1)$} & \textbf{Qubit dedotto} & \textbf{Esito} \\
\midrule
nessuno & $(0,0)$ & --- & \checkmark \\
su $q_0$ & $(1,0)$ & $q_0$ & \checkmark \\
su $q_1$ & $(1,1)$ & $q_1$ & \checkmark \\
su $q_2$ & $(0,1)$ & $q_2$ & \checkmark \\
\bottomrule
\end{tabular}
\caption[Verifica della sindrome del codice a ripetizione]{Lookup table
verificata per il codice bit-flip e per il duale phase-flip. L'ancilla
$a_0$ misura la parità di $\{0,1\}$ e $a_1$ quella di $\{1,2\}$. La
coincidenza cifra per cifra a parità di seed verifica l'implementazione
della dualità, non costituisce una conferma indipendente del fatto
algebrico.}
\label{tab:sindrome_repetition}
\end{table}

Il secondo controllo prepara lo stato logico, applica a ciascun qubit un
errore indipendente con probabilità $p$, estrae la sindrome, corregge e
decodifica. I $20\,000$ campioni per punto (seed 42) producono la curva
riportata nella Figura~\ref{fig:qec_repetition}. Un errore Z resta
invisibile al codice bit-flip, e viceversa per il duale: è il limite che
motiva il passaggio allo Steane.

\subsection{Struttura e protocollo completo del codice di Steane}
\label{app:qec_steane}

Lo Steane~\cite{steane1996} è un codice \ac{CSS} derivato dal codice di
Hamming $[7,4,3]$: codifica un qubit in sette, ha distanza 3 e corregge
un errore X, Z o Y su un singolo qubit.

\begin{table}[H]
\centering
\begin{tabular}{lcl}
\toprule
\textbf{Parametro} & \textbf{Valore} & \textbf{Significato} \\
\midrule
$n$ & 7 & qubit fisici \\
$k$ & 1 & qubit logici \\
$d$ & 3 & distanza \\
$t$ & 1 & errori arbitrari correggibili (Eq.~\ref{eq:distanza_t}) \\
Tipo & \ac{CSS} & sindromi X e Z separate \\
\bottomrule
\end{tabular}
\caption[Parametri del codice di Steane]{Parametri del codice di Steane $[[7,1,3]]$.}
\label{tab:steane_parametri}
\end{table}

Nella convenzione adottata i generatori sono
\begin{equation}
\begin{aligned}
\text{tipo X:}\quad & IIIXXXX,\quad IXXIIXX,\quad XIXIXIX,\\
\text{tipo Z:}\quad & IIIZZZZ,\quad IZZIIZZ,\quad ZIZIZIZ.
\end{aligned}
\label{eq:steane_stabilizzatori}
\end{equation}
Gli stabilizzatori Z rilevano gli errori X, quelli X rilevano gli errori
Z e un errore Y accende entrambe le famiglie. I tre bit di sindrome,
letti come numero binario, indicano la posizione del qubit colpito.

Il protocollo comprende: preparazione di $\ket{0_L}$ o $\ket{+_L}$;
iniezione manuale di X, Z e Y su ciascuno dei sette qubit; misura della
sindrome; costruzione della lookup table; stima della correzione al
variare di $p$; distinzione fra errori sui dati e sulle ancille. La
preparazione di $\ket{0_L}$ usa tre qubit seme in sovrapposizione e una
rete CNOT; i sei stabilizzatori restituiscono $000000$ su tutti i
$2\,000$ campioni, verificando l'appartenenza al code space.

\begin{table}[H]
\centering
\begin{tabular}{cccc}
\toprule
\textbf{Qubit} & \textbf{Sindrome} & \textbf{Valore} & \textbf{Correzione} \\
\midrule
$q_0$ & $001$ & 1 & applica $X$ (o $Z$) su $q_0$ \\
$q_1$ & $010$ & 2 & applica su $q_1$ \\
$q_2$ & $011$ & 3 & applica su $q_2$ \\
$q_3$ & $100$ & 4 & applica su $q_3$ \\
$q_4$ & $101$ & 5 & applica su $q_4$ \\
$q_5$ & $110$ & 6 & applica su $q_5$ \\
$q_6$ & $111$ & 7 & applica su $q_6$ \\
\bottomrule
\end{tabular}
\caption[Syndrome table del codice di Steane]{Tabella
sindrome--qubit--correzione verificata per tutti i 21 casi: X, Z e Y su
ciascuno dei sette qubit.}
\label{tab:steane_sindrome}
\end{table}

La simulazione Monte Carlo usa $200\,000$ campioni per punto (seed 42),
applica la correzione a peso minimo e verifica il residuo logico. La
versione implementata assume errori sui qubit dati e sindrome ideale:
l'estrazione non è fault-tolerant. Un guasto sulle ancille può corrompere
la sindrome e indurre una correzione errata; i cicli rumorosi e ripetuti
del surface code affrontano precisamente questo limite.

\end{onehalfspace}
